\documentclass[12pt,a4paper]{article}
\usepackage[margin=25mm, headheight=15pt, headsep=10pt]{geometry}
\usepackage{fontspec}
\usepackage[table]{xcolor} % Note: table option is required for rowcolors
\usepackage{titlesec}
\usepackage{tocloft}
\usepackage{fancyhdr}
\usepackage{booktabs}
\usepackage{tcolorbox}
\tcbuselibrary{skins, breakable}
\usepackage[colorlinks=true, linkcolor=emerald, urlcolor=emerald, bookmarks=true]{hyperref}

\definecolor{emerald}{RGB}{0,102,51}
\definecolor{darkred}{RGB}{139,0,0}
\definecolor{lightgray}{RGB}{245,245,245}

\usepackage[nil]{babel}
\babelprovide[import, main]{bengali}
\babelprovide[import]{arabic}
\babelfont{rm}[Renderer=HarfBuzz,Scale=1.0]{Noto Serif Bengali}
\babelfont{sf}[Renderer=HarfBuzz]{Noto Sans Bengali}
\babelfont{tt}[Renderer=HarfBuzz]{Noto Sans Bengali}
\babelfont[arabic]{rm}[Renderer=HarfBuzz,Scale=1.1]{Noto Naskh Arabic}

% Section styling
\titleformat{\section}{\Large\bfseries\color{emerald}}{\thesection.}{0.5em}{}
\titleformat{\subsection}{\large\bfseries\color{emerald!85!black}}{\thesubsection.}{0.5em}{}
\titleformat{\subsubsection}{\normalsize\bfseries\color{emerald!70!black}}{\thesubsubsection.}{0.5em}{}
\titlespacing*{\section}{0pt}{18pt}{8pt}

% Fancy Header/Footer setup
\pagestyle{fancy}
\fancyhf{} % clear all header and footer fields
\fancyhead[L]{\color{emerald}\small\textbf{\nouppercase{\leftmark}}}
\fancyfoot[L]{\scriptsize\color{black!50}{\lat{AI}}-সংকলিত, আলেম-যাচাইকৃত নয়}
\fancyfoot[C]{\color{emerald!70!black}\thepage}
\renewcommand{\headrulewidth}{0.4pt}
\renewcommand{\headrule}{\hbox to\headwidth{\color{emerald}\leaders\hrule height \headrulewidth\hfill}}
\renewcommand{\footrulewidth}{0pt}

% Custom boxes
% 1. Ayat/Hadith Box
\newtcolorbox{ayatbox}[1][]{
    enhanced, breakable, 
    colback=emerald!5, 
    colframe=emerald, 
    boxrule=0.5pt, 
    arc=3pt, 
    left=10pt, right=10pt, top=10pt, bottom=10pt,
    #1
}

% 2. Quote/Translation Box
\newtcolorbox{quotebox}[1][]{
    enhanced, breakable,
    frame hidden,
    colback=lightgray,
    borderline west={3pt}{0pt}{emerald},
    left=15pt, right=10pt, top=10pt, bottom=10pt,
    sharp corners,
    #1
}

% 3. AI-generated notice box
\newtcolorbox{warnbox}[1][]{
    enhanced, breakable,
    colback=red!4,
    colframe=darkred,
    boxrule=0.8pt,
    arc=3pt,
    left=10pt, right=10pt, top=10pt, bottom=10pt,
    #1
}

% Commands
\newcommand{\arab}[1]{\foreignlanguage{arabic}{#1}}
\newcommand{\kib}[1]{\textcolor{emerald}{\textbf{#1}}}

% Latin (English) fallback: Noto Serif Bengali has NO Latin glyphs
\newfontfamily\latfont[Renderer=HarfBuzz]{Noto Serif}
\newcommand{\lat}[1]{{\latfont #1}}

\setlength{\parskip}{5pt}
\setlength{\parindent}{0pt}

% Fix for missing bullet in Noto Serif Bengali
\renewcommand{\labelitemi}{$\bullet$}



\begin{document}
\begin{center}{\Large\bfseries\color{emerald} \lat{04}-সূরা-ফাতিহা-অর্থ}\end{center}
\vspace{1em}
\section*{০৪ — সূরা ফাতিহা: \textit{\lat{Al-Fatihah}} — নামাজের রুকন (প্রতিটি রাকাতে পঠিত)}

\begin{warnbox}
 \textbf{গুরুত্বপূর্ণ নোটিশ:} এই কন্টেন্টটি \lat{AI} দিয়ে প্রস্তুত। কেবল সহীহ/হাসান হাদিস রাখা হয়েছে। আলেম-যাচাইকৃত নয়।
\end{warnbox}

\medskip\hrule\medskip

\subsection*{১. সূত্র কার্ড}

\begin{table}[h]\centering\small
\begin{tabular}{ll}
বিষয় & বিবরণ \\
\hline
\textbf{সূরা} & আল-ফাতিহা (১) — মক্কী \\
\textbf{আয়াত} & ৭ আয়াত \\
\textbf{পারা} & ১ \\
\textbf{নাম} & উম্মুল কিতাব, সাবা'ুল মাসানী, আল-হামদ, আশ-শিফা \\
\textbf{বৈশিষ্ট্য} & \textbf{নামাজের রুকন} — ছাড়লে নামাজ নয় \\
\textbf{সূত্র} & সহীহ বুখারি ৭৫৬, সহীহ মুসলিম ৩৯৪ \\
\hline
\end{tabular}
\end{table}

\medskip\hrule\medskip

\subsection*{২. মূল আরবি}

\begin{center}\arab{بِسْمِ اللَّهِ الرَّحْمَٰنِ الرَّحِيمِ}\end{center}
\begin{center}\arab{الْحَمْدُ لِلَّهِ رَبِّ الْعَالَمِينَ}\end{center}
\begin{center}\arab{الرَّحْمَٰنِ الرَّحِيمِ}\end{center}
\begin{center}\arab{مَالِكِ يَوْمِ الدِّينِ}\end{center}
\textbf{\arab{إِيَّاكَ} \arab{نَعْبُدُ} \arab{وَإِيَّاكَ} \arab{نَ}�\arab{ْتَعِينُ}}
\begin{center}\arab{اهْدِنَا الصِّرَاطَ الْمُسْتَقِيمَ}\end{center}
\begin{center}\arab{صِرَاطَ الَّذِينَ أَنْعَمْتَ عَلَيْهِمْ غَيْرِ الْمَغْضُوبِ عَلَيْهِمْ وَلَا الضَّالِّينَ}\end{center}

\subsection*{৩. বাংলা উচ্চারণ}

\textbf{বিসমিল্লাহির রাহমানির রাহীম}
\textbf{আলহামদু লিল্লাহি রাব্বিল আলামীন}
\textbf{আর রাহমানির রাহীম}
\textbf{মালিকি যাওমিদ দীন}
\textbf{ইয়াকা নাআবুদু ওয়া ইয়াকা নাস্তাঈন}
\textbf{ইহদিনাস সিরাতাল মুস্তাকীম}
\textbf{সিরাতাল্লাজীনা আনআমতা আলাইহিম গৈরিল মাগদুবি আলাইহিম ওয়ালাদ্বাল্লীন}

\medskip\hrule\medskip

\subsection*{৪. শব্দ-শব্দ অর্থ}

\begin{table}[h]\centering\small
\begin{tabular}{lll}
আরবি & বাংলা & \lat{English} \\
\hline
\textbf{\arab{بِسْمِ}} & নামে & \textit{\lat{in} \lat{the} \lat{name}} \\
\textbf{\arab{اللَّهِ}} & আল্লাহর & \textit{\lat{of} \lat{Allah}} \\
\textbf{\arab{الرَّحْمَٰنِ}} & পরম করুণাময় & \textit{\lat{the} \lat{Most} \lat{Gracious}} \\
\textbf{\arab{الرَّحِيمِ}} & পরম দয়ালু & \textit{\lat{the} \lat{Most} \lat{Merciful}} \\
\textbf{\arab{الْحَمْدُ}} & সকল প্রশংসা & \textit{\lat{all} \lat{praise}} \\
\textbf{\arab{لِلَّهِ}} & আল্লাহর জন্যে & \textit{\lat{is} \lat{for} \lat{Allah}} \\
\textbf{\arab{رَبِّ}} & পালনকর্তা & \textit{\lat{Lord}} \\
\textbf{\arab{الْعَالَمِينَ}} & সকল জগত & \textit{\lat{of} \lat{all} \lat{worlds}} \\
\textbf{\arab{مَالِكِ}} & মালিক/স্বামী & \textit{\lat{Master}} \\
\textbf{\arab{يَوْمِ} \arab{الدِّينِ}} & বিচারদিন & \textit{\lat{Day} \lat{of} \lat{Judgement}} \\
\textbf{\arab{إِيَّاكَ}} & শুধু তোমাকেই & \textit{\lat{You} \lat{alone}} \\
\textbf{\arab{نَعْبُدُ}} & আমরা ইবাদত করি & \textit{\lat{we} \lat{worship}} \\
\textbf{\arab{نَسْتَعِينُ}} & আমরা সাহায্য চাই & \textit{\lat{we} \lat{seek} \lat{help}} \\
\textbf{\arab{اهْدِ}ন\arab{َا}} & পথ দেখাও & \textit{\lat{guide} \lat{us}} \\
\textbf{\arab{الصِّرَاطَ} \arab{الْمُسْتَقِيمَ}} & সরল পথটি & \textit{\lat{the} \lat{straight} \lat{path}} \\
\textbf{\arab{أَنْعَمْتَ} \arab{عَلَيْهِمْ}} & যাদের ওপর অনুগ্রহ করেছ & \textit{\lat{those} \lat{upon} \lat{whom} \lat{You} \lat{bestowed} \lat{favour}} \\
\textbf{\arab{الْمَغْضُوبِ}} & কূপিতদের & \textit{\lat{those} \lat{who} \lat{earned} \lat{anger}} \\
\textbf{\arab{الضَّالِّينَ}} & পথভ্রষ্টদের & \textit{\lat{those} \lat{who} \lat{went} \lat{astray}} \\
\hline
\end{tabular}
\end{table}

\medskip\hrule\medskip

\subsection*{৫. আয়াত-বার পূর্ণ বাংলা অনুবাদ}

\textbf{১। \lat{“}আল্লাহর নামে, যিনি পরম করুণাময়, পরম দয়ালু।\lat{”}}

\textbf{২। \lat{“}সকল প্রশংসা আল্লাহর জন্যে — যিনি সকল জগতের পালনকর্তা।\lat{”}}

\textbf{৩। \lat{“}পরম করুণাময়, পরম দয়ালু।\lat{”}}

\textbf{৪। \lat{“}বিচারদিনের মালিক।\lat{”}}

\textbf{৫। \lat{“}শুধু তোমাকেই আমরা ইবাদত করি এবং শুধু তোমাকেই আমরা সাহায্য চাই।\lat{”}}

\textbf{৬। \lat{“}আমাদের সরল পথে হেদায়াত কর।\lat{”}}

\textbf{৭। \lat{“}যাদের ওপর তুমি অনুগ্রহ করেছ — তাদের পথ, যারা কূপিত হয়েছ তাদের নয় এবং পথভ্রষ্টদেরও নয়।\lat{”}}

\medskip\hrule\medskip

\subsection*{৬. সংক্ষিপ্ত ব্যাখ্যা}

\subsubsection*{\textbf{আয়াত ১-২: হামদ (প্রশংসা)}}
নামাজ শুরু হচ্ছে আল্লাহর প্রশংসা দিয়ে। \textbf{\lat{Rabb} (রব)} মানে যিনি সৃষ্টি, পোষণ, পরিচালনা করেন। "\lat{Al-alamin}" — সকল জগত: মানুষ, জিন, ফেরেশতা, প্রাণী, গ্রহ-নক্ষত্র।

\subsubsection*{\textbf{আয়াত ৩: রাহমান ও রাহীম}}
\textbf{\lat{Ar-Rahman}} — সকল সৃষ্টির জন্য সাধারণ দয়া। \textbf{\lat{Ar-Rahim}} — মুমিনদের জন্য বিশেষ দয়া (আখেরাতে)।

\subsubsection*{\textbf{আয়াত ৪: মালিকি যাওমিদ দীন}}
"\lat{Malik}" — মালিক, বিচারক। "\lat{Yawmid} \lat{Din}" — বিচারদিন। এতে মানুষের মনে \textbf{\lat{responsibility} (দায়িত্ব)} জাগে।

\subsubsection*{\textbf{আয়াত ৫: ইয়াকা নাআবুদু ওয়া ইয়াকা নাস্তাঈন}}
\textbf{"শুধু তোমাকেই ইবাদত করি"} — \lat{Tawhid} (তাওহীদ), \lat{ikhlas} (ইখলাস)। \textbf{"শুধু তোমাকেই সাহায্য চাই"} — \lat{Tawakkul} (তাওয়াক্কুল, আল্লাহর ওপর ভরসা)।

\subsubsection*{\textbf{আয়াত ৬: ইহদিনাস সিরাতাল মুস্তাকীম}}
"\lat{Sirat}" — পথ। "\lat{Mustaqim}" — সীধা, বাঁকা নয়। প্রতিটি রাকাতে এই দোয়া — কারণ \lat{guidance} (হেদায়াত) প্রতিটি মুহূর্তে প্রয়োজন।

\subsubsection*{\textbf{আয়াত ৭: তিন শ্রেণী মানুষ}}
১। \textbf{আনামতা আলাইহিম} — যাদের ওপর অনুগ্রহ (নবী, সিদ্দিক, শহীদ, সালেহ — সূরা আন-নিসা ৪:৬৯)।
২। \textbf{মাগদুবি আলাইহিম} — যারা জানে কিন্তু মানে না (\lat{knowledge} \lat{without} \lat{action})।
৩। \textbf{আদ-দাল্লীন} — যারা জানে না, পথভ্রষ্ট (\lat{ignorance})।

\medskip\hrule\medskip

\subsection*{৭. খুশুর জন্য মূল পয়েন্টস}

\begin{table}[h]\centering\small
\begin{tabular}{ll}
পয়েন্ট & কীভাবে প্রয়োগ করবেন \\
\hline
\textbf{বিসমিল্লাহ} & ধীরে পড়ুন — "আমি আল্লাহর নামে শুরু করছি" \\
\textbf{আলহামদু} & হৃদয়ে \lat{gratitude} (কৃতজ্ঞতা) আনুন \\
\textbf{ইয়াকা নাআবুদু} & মনে করুন: "আমার কোনো অন্য নেই, শুধু আল্লাহ" \\
\textbf{ইহদিনা} & চোখ বন্ধ করে দোয়া করুন: "হে আল্লাহ, পথ থেকে ভটকে যাওয়া থেকে বাঁচাও" \\
\textbf{আমিন} & শেষে নীরবে "আমিন" বলুন — ফেরেশতাও বলে (বুখারি ৭৮০) \\
\hline
\end{tabular}
\end{table}

\medskip\hrule\medskip

\subsection*{৮. সংশ্লিষ্ট হাদিস}

\begin{table}[h]\centering\small
\begin{tabular}{ll}
হাদিস & গ্রন্থ \\
\hline
\textbf{\lat{“}যে ব্যক্তি সূরা ফাতিহা পড়েনি, তার নামাজ নেই।\lat{”}} & সহীহ বুখারি ৭৫৬ \\
\textbf{\lat{“}নামাজ আমি ও বান্দার মাঝে আধা-আধা বাটলাম...\lat{”}} (হাদিসে কুদসি) & সহীহ মুসলিম ৩৯৫ \\
\textbf{\lat{“}যখন ইমাম 'ওয়ালাদ্দাল্লিন' বলেন, তোমরা আমিন বলো...\lat{”}} & সহীহ বুখারি ৭৮০ \\
\textbf{\lat{“}সূরা ফাতিহা সব রোগের চিকিৎসা।\lat{”}} & সহীহ বুখারি ৫৭৩৬ \\
\hline
\end{tabular}
\end{table}

\medskip\hrule\medskip

\textit{শেষ আপডেট: আগস্ট ২০২৬}


\end{document}
